\begin{helper}
Let $\mathcal F$ be a $k$-uniform hypergraph and let $f\in E(\mathcal F)$.
Suppose that every vertex of $f$ has degree exactly $D$, and that the number
of line-graph neighbors of $f$ is exactly $kD-k$.  Then every edge
$g\ne f$ meeting $f$ meets it in exactly one vertex:
\[
|g\cap f|=1 .
\]
\end{helper}
\begin{proof}
Let $N$ be the set of edges $g\ne f$ with $g\cap f\ne\emptyset$.  For a
vertex $v\in f$, the number of edges other than $f$ containing $v$ is
$D-1$, because $v$ has degree $D$ and one of those incident edges is $f$.
Since $\mathcal F$ is $k$-uniform, $|f|=k$, and hence the total number of
incidences
\[
\{(v,g): v\in f,\ g\in N,\ v\in g\}
\]
is $k(D-1)$.

Counting the same incidences by the edge $g\in N$ gives
\[
\sum_{g\in N}|g\cap f|=k(D-1).
\]
The hypothesis on the line-neighborhood says
$|N|=kD-k=k(D-1)$.  Each summand $|g\cap f|$ is at least $1$, by the
definition of $N$.  If some $g\in N$ had $|g\cap f|\ne 1$, then since it is
a positive integer it would have $|g\cap f|\ge2$, while all other summands
would still be at least $1$.  The sum would then be at least $|N|+1$,
contradicting $\sum_{g\in N}|g\cap f|=|N|$.  Therefore every $g\in N$
satisfies $|g\cap f|=1$.
\end{proof}
